\documentclass[10pt,a4paper]{amsart}
\usepackage[margin=19mm]{geometry}
\usepackage{amsmath,amssymb,mathtools}
\usepackage[scr=rsfs,cal=euler]{mathalfa}
\usepackage{xfrac}
\usepackage[unicode,hidelinks,bookmarks=false]{hyperref}
\newcommand{\ie}{i.e.\ }
\newcommand{\cf}{cf.\ }
\newcommand{\resp}{respectively\ }
\newcommand{\Subsection}{\subsection}
\providecommand{\nobreakdash}{\nobreak\hbox{-}}
\setlength{\emergencystretch}{2em}
\multlinegap0pt
\usepackage{amssymb}
\newtheorem{theorem}{Theorem}
\title{full title}
\author{Zhi-Lin Dai, Hai-Ping Fu$^*$, Yao Lu}
\date{Source: arXiv:2609.05106v1 (CC BY 4.0)}
\begin{document}
\maketitle

\noindent\emph{Source and licence.} full title. Version arXiv:2609.05106v1 (CC BY 4.0), distributed by arXiv under the Creative Commons Attribution 4.0 International licence, http://creativecommons.org/licenses/by/4.0/, as recorded in the arXiv licence metadata for this exact version and shown on the abstract page https://arxiv.org/abs/2609.05106v1. Full licence notice: \texttt{licenses/arxiv-2609.05106-CC-BY-4.0.txt}.\par\smallskip

\noindent\textit{Editorial scope.} Counted unit: the article's theorem theorem1.1 (source0.tex lines 141--153). The article's complete original proof of it is delivered with it (source0.tex lines 348--431, one \texttt{proof} environment of 84 lines, which the article places away from the statement). No mathematics is written, repaired or re-derived: the statement and the proof are the article's own lines, in the article's own notation, and no retained line is substituted or re-flowed.  No further source-local context is needed: every cross-reference of every retained line resolves inside the two delivered spans.  Nothing was omitted from any retained span, so the package carries no omission record.  No retained line contains a citation, so the wrapper carries no bibliography and no bibliography entry.  No retained line contains a figure environment, an image-inclusion command or drawing code, so no image is delivered and the package declares no asset dependency.  Editorial formatting only, in addition: an AMS \texttt{amsart} preamble that restates the article's own declarations and macros used by the retained lines, namely the environments \texttt{theorem} on separate counters, together with the standard packages those lines need.  The retained environments print in this document as Theorem 1 in order of appearance, whereas the article prints the same items with the numbering of its own full document.  The complete native source of the article as distributed in the arXiv e-print for this exact version is bundled unchanged as \texttt{sources/209415/source0.tex}.
\begin{theorem}
\label{theorem1.1}
Let $\left( {M,g,J} \right)$ be a K{\"a}hler manifold of complex dimension $n$.
\begin{enumerate}
\item If $M$ has $\frac{{n\left(n + 1\right)}}{2}$-positive (resp., nonnegative) Calabi curvature operator, then $M$ has positive (resp., nonnegative) scalar curvature.
\item If $M$ has $\frac{{n + 1}}{2}$-positive (resp., nonnegative) Calabi curvature operator, then $M$ has positive (resp., nonnegative) Ricci curvature.
\item If $M$ has $\left(n-1\right)$-positive (resp., nonnegative) Calabi curvature operator, then $M$ has positive (resp., nonnegative) orthogonal Ricci curvature.
\item If $M$ has $n$-positive (resp., nonnegative) Calabi curvature operator, then at every point it satisfies
$$2{Ric\left( {{e_b},{e_b}} \right) - R\left( {{e_b},J{e_b},{e_b},J{e_b}} \right)}> \left( \ge \right) 0.$$

\item If $M$ has $\alpha $-positive Calabi curvature operator for some $\alpha  \le n-1$, then $M$ is irreducible.
\end{enumerate}
\end{theorem}

\begin{proof}[{\bf Proof of Theorem~\ref{theorem1.1}}]
(1) Take the orthonormal basis of ${ \odot ^2}{V^{1,0}}$ given by $${\left\{ {\frac{1}{{\sqrt 2 }}\left( {{Z_A} \odot {Z_B}} \right)} \right\}_{1 \le A < B \le n}} \cup \left\{ {\frac{1}{2}\left( {{Z_A} \odot {Z_A}} \right)} \right\}_{A = 1}^n,$$ then

\begin{align*}
  &\sum\limits_{1 \le A < B \le n}^{} {g\left( {\mathcal{C} \left( {\frac{1}{{\sqrt 2 }}\left( {{Z_A} \odot {Z_B}} \right)} \right),\frac{1}{{\sqrt 2 }}\left( {\overline {{Z_A}}  \odot \overline {{Z_B}} } \right)} \right)}  + \sum\limits_{A = 1}^n {g\left( {\mathcal{C}\left( {\frac{1}{2}\left( {{Z_A} \odot {Z_A}} \right)} \right),\frac{1}{2}\left( {\overline {{Z_A}}  \odot \overline {{Z_A}} } \right)} \right)}  \\
  &= \frac{1}{4}\sum\limits_{A,B = 1}^n {g\left( {\mathcal{C}\left( {{Z_A} \odot {Z_B}} \right),\overline {{Z_A}}  \odot \overline {{Z_B}} } \right)}  =\sum\limits_{A,B = 1}^n {R\left( {{Z_A},\overline {{Z_A}} ,\overline {{Z_B}} ,{Z_B}} \right)}\\
   &=\frac{1}{2}\sum\limits_{a,B = 1}^n {R\left( {{e_a} - \sqrt { - 1} J{e_a},{e_a} + \sqrt { - 1} J{e_a},\overline {{Z_B}} ,{Z_B}} \right)}
   = \sum\limits_{a,B = 1}^n {R\left( {{e_a},\sqrt { - 1} J{e_a},\overline {{Z_B}} ,{Z_B}} \right)} \\
   &= \sum\limits_{a,b = 1}^n {R\left( {{e_a},J{e_a},{e_b},J{e_b}} \right)}
   = \sum\limits_{a = 1}^n {Ric\left( {{e_a},{e_a}} \right)}
   =\frac{s}{2} ,
\end{align*}
because $M$ has $\frac{{n\left( {n + 1} \right)}}{2}$-positive ( nonnegative) Calabi curvature operator,  the scalar curvature is positive(nonnegative).

(2) Let $A \in \left\{ {1,2, \cdots ,n} \right\}$ be fixed. For all $B \ne A$, define
$$2{\mu _b}=g\left( {\mathcal{C}\left( {\frac{1}{{\sqrt 2 }}\left( {{Z_A} \odot {Z_B}} \right)} \right),\frac{1}{{\sqrt 2 }}\left( {\overline {{Z_A}}  \odot \overline {{Z_B}} } \right)} \right) = 2R\left( {{e_a},J{e_a},{e_b},J{e_b}} \right),$$
$${\mu _a}=g\left( {\mathcal{C}\left( {\frac{1}{2}\left( {{Z_A} \odot {Z_A}} \right)} \right),\frac{1}{2}\left( {\overline {{Z_A}}  \odot \overline {{Z_A}} } \right)} \right) = R\left( {{e_a},J{e_a},{e_a},J{e_a}} \right).$$
Arranging these $n$ numbers in ascending order yields
$$2{\mu _1}, \cdots 2{\mu _{k - 1}},{\mu _k} = {\mu _a},2{\mu _{k + 1}}, \cdots, 2{\mu _n}.$$
Because $M$ has $\frac{n+1}{2}$- positive(nonnegative) Calabi curvature operator, we obtain the following estimates.

Case 1. $k\le \frac{n+1}{2}$.  When $n$ is odd,
\begin{align*}
  0 \le & \sum\limits_{i = 1}^{k - 1} {2\mu _i}  + {\mu _k} + \sum\limits_{j = k + 1}^{\frac{{n + 1}}{2}} {2{\mu _j}} \\
   \le & \sum\limits_{i = 1}^{k - 1} {\left( {{\mu _i} + {\mu _{n - i + 1}}} \right)}  + {\mu _k} + \sum\limits_{j = k + 1}^{\frac{{n + 1}}{2}} {\left( {{\mu _j} + {\mu _{n - j + 2}}} \right)} \\
   =& \sum\limits_{l = 1}^n {{\mu _l}}  = \sum\limits_{b = 1}^n {R\left( {{e_a},J{e_a},{e_b},J{e_b}} \right)}
   =Ric\left( {{e_a},{e_a}} \right).
\end{align*}
When $n$ is even,
\begin{align*}
  0 \le & \sum\limits_{i = 1}^{k - 1} {2\mu _i}  + {\mu _k} + \sum\limits_{j = k + 1}^{\frac{n}{2}} {2{\mu _j}}  + {\mu _{\frac{n}{2} + 1}}  \\
   \le & \sum\limits_{i = 1}^{k - 1} {\left( {{\mu _i} + {\mu _{n - i + 1}}} \right)}  + {\mu _k} + \sum\limits_{j = k + 1}^{\frac{n}{2}} {\left( {{\mu _j} + {\mu _{n - j + 2}}} \right)}  + {\mu _{\frac{n}{2} + 1}} \\
   =& \sum\limits_{l = 1}^n {{\mu _l}}  = \sum\limits_{b = 1}^n {R\left( {{e_a},J{e_a},{e_b},J{e_b}} \right)}
   = Ric\left( {{e_a},{e_a}} \right) .
\end{align*}
Case 2.  $k > \frac{n+1}{2}$. When $n$ is odd,
\begin{align*}
0 \le & \sum\limits_{i = 1}^{\frac{{n + 1}}{2}} {2{\mu _i}}
 \le  \sum\limits_{i = 1}^{n - k} {\left( {{\mu _i} + {\mu _{n - i + 1}}} \right)}  + \left( {{\mu _{n - k + 1}} + {\mu _{n - k + 2}} + {\mu _k}} \right) + \sum\limits_{j = n - k + 3}^{\frac{{n + 1}}{2}} {\left( {{\mu _j} + {\mu _{n - j + 2}}} \right)} \\
 = & \sum\limits_{l = 1}^n {{\mu _l}}  = \sum\limits_{b = 1}^n {R\left( {{e_a},J{e_a},{e_b},J{e_b}} \right)}
 =  Ric\left( {{e_a},{e_a}} \right).
\end{align*}
When $n$ is even and $k \ne \frac{n}{2}+1$,
\begin{align*}
0 \le & \sum\limits_{i = 1}^{\frac{n}{2}} {2{\mu _i}}  + {\mu _{\frac{n}{2} + 1}}\\
 \le & \sum\limits_{i = 1}^{n - k} {\left( {{\mu _i} + {\mu _{n - i + 1}}} \right)}  + \left( {{\mu _{n - k + 1}} + {\mu _{n - k + 2}} + {\mu _k}} \right) + \sum\limits_{j = n - k + 3}^{\frac{n}{2}} {\left( {{\mu _j} + {\mu _{n - j + 2}}} \right)}  + {\mu _{\frac{n}{2} + 1}}\\
 = & \sum\limits_{l = 1}^n {{\mu _l}}  = \sum\limits_{b = 1}^n {R\left( {{e_a},J{e_a},{e_b},J{e_b}} \right)}
 =  Ric\left( {{e_a},{e_a}} \right).
\end{align*}
When $n$ is even and $k = \frac{n}{2}+1$,
\begin{align*}
0 \le & \sum\limits_{i = 1}^{\frac{n}{2}} {2{\mu _i}}  + \frac{1}{2}{\mu _{\frac{n}{2} + 1}}\\
 \le & \sum\limits_{i = 1}^{\frac{n}{2} - 1} {\left( {{\mu _i} + {\mu _{n - i + 1}}} \right)}  + \left( {{\mu _{\frac{n}{2}}} + \frac{1}{2}{\mu _{\frac{n}{2} + 1}}} \right) + \frac{1}{2}{\mu _{\frac{n}{2} + 1}}\\
 = & \sum\limits_{l = 1}^n {{\mu _l}}  = \sum\limits_{b = 1}^n {R\left( {{e_a},J{e_a},{e_b},J{e_b}} \right)}
 =  Ric\left( {{e_a},{e_a}} \right).
\end{align*}
Combining all cases and using   $Ric\left( {J{e_a},J{e_a}} \right) = Ric\left( {{e_a},{e_a}} \right)$, we can conclude that $Ric \ge 0$.

(3) Since $M$ has $\left( n-1 \right)$-positive(nonnegative) Calabi curvature operator, we obtain
\begin{align*}
0 \le & \sum\limits_{b \ne a}^{} {2{\mu _b}}
 = \sum\limits_{b = 1}^n {2{\mu _b}}  - 2{\mu _a}\\
 =& 2\sum\limits_{b = 1}^n {R\left( {{e_a},J{e_a},{e_b},J{e_b}} \right)}  - 2R\left( {{e_a},J{e_a},{e_a},J{e_a}} \right)\\
 =& 2\left( {Ric\left( {{e_a},{e_a}} \right) - R\left( {{e_a},J{e_a},{e_a},J{e_a}} \right)} \right)
 = 2Ri{c^ \bot }\left( {{e_a},{e_a}} \right) .
\end{align*}

(4) Since $M$ has $ n $-positive(nonnegative) Calabi curvature operator, we obtain
\begin{align*}
0 \le & \sum\limits_{b \ne a}^{} {2{\mu _b}}  + {\mu _a}
 =  \sum\limits_{b = 1}^n {2{\mu _b}}  - {\mu _a}\\
 = & 2\sum\limits_{b = 1}^n {R\left( {{e_a},J{e_a},{e_b},J{e_b}} \right)}  - R\left( {{e_a},J{e_a},{e_a},J{e_a}} \right)\\
 = & 2Ric\left( {{e_a},{e_a}} \right) - R\left( {{e_a},J{e_a},{e_a},J{e_a}} \right).
\end{align*}


(5) Suppose $M$ is reducible, and write $M = M_1 \times M_2$, where $M_i$ is a complex $n_i$-dimensional K{\"a}hler manifold, $i=1,2$. Choose an orthonormal basis $$\left\{ {{e_1}, \cdots ,{e_{{n_1}}}} , {{Je_1}, \cdots ,{Je_{{n_1}}}} \right\}$$ for $M_1$ and $$\left\{ {{e_{{n_1} + 1}}, \cdots ,{e_{{n_1} + {n_2}}}},{{Je_{{n_1} + 1}}, \cdots ,{Je_{{n_1} + {n_2}}}} \right\}$$ for $ M_2$. For indices $1 \le A \le {n_1}$ and ${n_1} + 1 \le C \le {n_1} + {n_2}$, the curvature tensor of a product satisfies
$$ R\left( {{Z_A},\overline {{Z_A}} ,\overline {{Z_C}} ,{Z_C}} \right) = 0, $$
hence
$$g\left( {\mathcal{C}\left( {\frac{1}{{\sqrt 2 }}{Z_A} \odot {Z_C}} \right),\frac{1}{{\sqrt 2 }}\overline {{Z_A}}  \odot \overline {{Z_C}} } \right) =0.$$
Consequently,
$$\sum\limits_{\alpha  = 1}^{{n_1}  {n_2}} {{\lambda _\alpha }}  \le \sum\limits_{A = 1}^{{n_1}} {\sum\limits_{C = {n_1} + 1}^{{n_1} + {n_2}} {g\left( {\mathcal{C}\left( {\frac{1}{{\sqrt 2 }}{Z_A} \odot {Z_C}} \right),\frac{1}{{\sqrt 2 }}\overline {{Z_A}}  \odot \overline {{Z_C}} } \right)} }  = 0.$$
Thus the Calabi operator  $\mathcal{C}$ is $ {n_1}{n_2} $-nonpositive . This contradicts the hypothesis that $${\lambda _1} + {\lambda _2} +  \cdots  + {\lambda _\alpha} > 0\ \ \  \text{for}\ \   \alpha \le n-1.$$ Therefore $M$ must be irreducible.
\end{proof}

\end{document}
